\section{Bayesian Analysis}

%--------------------------------------------------------
\begin{frame}
\frametitle{Bayesian Parameter Estimation}

\small

The properties of the QGP are not directly observable; they are inferred by
comparing theoretical predictions with experimental measurements.

\vspace{0.15cm}

\begin{columns}[T]

%--------------------------------------------------------
\column{0.50\textwidth}

Bayes' theorem updates prior knowledge using experimental evidence:

\[
P(\boldsymbol{\theta}|D)
\propto
P(D|\boldsymbol{\theta})P(\boldsymbol{\theta})
\]

\begin{itemize}
    \item $\boldsymbol{\theta}$: model parameters
    \item $P(\boldsymbol{\theta})$: prior
    \item $P(D|\boldsymbol{\theta})$: likelihood
    \item $P(\boldsymbol{\theta}|D)$: posterior
\end{itemize}

{\tiny
J. E. Bernhard, J. S. Moreland and S. A. Bass,
\textit{Bayesian estimation of the specific shear and bulk viscosity of quark--gluon plasma},
Nature Physics \textbf{15}, 1113--1117 (2019).
}
%--------------------------------------------------------
\column{0.50\textwidth}

\centering
\footnotesize

\[
\boxed{\text{Experimental Data}}
\hspace{0.6cm}
\boxed{\text{Model Predictions}}
\]

\[
\searrow
\hspace{2.0cm}
\swarrow
\]

\[
\boxed{\text{Likelihood}}
\]

\[
\Downarrow
\]

\[
\boxed{\text{Posterior}}
\]

Sampled through MCMC.

\end{columns}

\vfill

\end{frame}